\section{Materiales y Métodos}
\label{sec:metodos}

\subsection{Design Science Research instanciado}

Este proyecto sigue el modelo de proceso de seis actividades de Peffers et al. \cite{peffers2007}. La
Tabla~\ref{tab:dsr} instancia cada actividad contra lo que realmente se hizo, con evidencia verificable.

\begin{table}[H]
\centering
\small
\begin{tabular}{p{2.6cm}p{9.5cm}}
\toprule
\textbf{Actividad DSR} & \textbf{Instanciación real en este proyecto} \\
\midrule
1. Identificar el problema y motivar & Proceso manual de pre-sustentación sin trazabilidad (Sección~\ref{sec:introduccion}); motivado por observaciones docentes reales de Entregas 1A/1B (\texttt{docs/observaciones/OBSERVACIONES.md}, 7 hallazgos con hash de commit) \\
2. Definir los objetivos de una solución & 12 requisitos funcionales priorizados MoSCoW (Sección~\ref{sec:requisitos}) \\
3. Diseño y desarrollo & Arquitectura de 3 capas + 7 ADR (Sección~\ref{sec:arquitectura}); implementación real (Sección~\ref{sec:implementacion}) \\
4. Demostración & Sistema funcional verificado end-to-end (login real, flujo de solicitud-jurado-cronograma-evaluación-acta) contra la topología de producción (nginx + backend + Postgres + Redis) \\
5. Evaluación & Evaluación empírica multi-dimensional: rendimiento (k6), seguridad (OWASP ZAP + find-sec-bugs), cobertura (JaCoCo), calidad web (Lighthouse) --- Sección~\ref{sec:evaluacion} \\
6. Comunicación & Este informe, el SRS v1.0.0, los ADR, y el repositorio público con licencia MIT \\
\bottomrule
\end{tabular}
\caption{Instanciación del modelo de proceso DSR de Peffers et al. \cite{peffers2007}.}
\label{tab:dsr}
\end{table}

\textbf{Nota honesta sobre las iteraciones DSR:} el modelo de Peffers es explícitamente iterativo (el
resultado de la evaluación puede regresar a cualquier actividad anterior). Este proyecto sí iteró en la
práctica --- por ejemplo, la evaluación de seguridad de la Fase 5 (actividad 5) encontró un secreto JWT
hardcodeado que obligó a regresar a la actividad 3 (rediseño de la configuración de seguridad) --- pero
esas iteraciones no se planificaron como tales desde el inicio del proyecto; se documentan aquí de forma
retrospectiva, con la fecha y el hallazgo real que las disparó, no como un proceso DSR diseñado
formalmente desde la Entrega 1A.

\subsection{Esquema Goal-Question-Metric}

Se instancia un esquema GQM \cite{basili1994} completo para la meta de rendimiento bajo carga, la única
dimensión de la evaluación empírica con datos suficientes para un análisis estadístico inferencial real
(Tabla~\ref{tab:gqm}).

\begin{table}[H]
\centering
\small
\begin{tabular}{p{2cm}p{10cm}}
\toprule
\textbf{Nivel} & \textbf{Contenido} \\
\midrule
\textbf{Meta (Goal)} & Analizar el sistema de caché Redis con el propósito de caracterizar su efecto
sobre la latencia, respecto al tiempo de respuesta del endpoint \texttt{/api/v1/universidades}, desde el
punto de vista del equipo de desarrollo, en el contexto de carga moderada (50 usuarios concurrentes) \\
\textbf{Pregunta 1 (Question)} & ¿Cuál es la latencia típica (mediana, percentiles) del endpoint en
estado de caché fría vs. caliente? \\
\textbf{Métrica 1.1} & Percentiles p50/p90/p95/p99 de latencia (ms), $n=30$ por escenario \\
\textbf{Métrica 1.2} & Media y desviación estándar, con intervalo de confianza del 95\,\% \\
\textbf{Pregunta 2 (Question)} & ¿La diferencia observada entre caché fría y caliente es
estadísticamente significativa, y de qué magnitud práctica? \\
\textbf{Métrica 2.1} & Estadístico U y valor p del test de Mann-Whitney/Wilcoxon (dos muestras
independientes) \\
\textbf{Métrica 2.2} & Tamaño de efecto (correlación biserial de rangos) \\
\textbf{Pregunta 3 (Question)} & ¿El sistema cumple el requisito no funcional RNF-01 (respuesta $<$
500\,ms bajo 50 usuarios concurrentes)? \\
\textbf{Métrica 3.1} & p95 de \texttt{http\_req\_duration} en cada una de las 5 corridas k6, comparado
contra el umbral de 500\,ms \\
\bottomrule
\end{tabular}
\caption{Esquema GQM para la meta de rendimiento del sistema de caché.}
\label{tab:gqm}
\end{table}

Los resultados de este esquema GQM se presentan en la Sección~\ref{sec:evaluacion} y se reproducen
ejecutando \texttt{scripts/perf-analysis.ipynb} contra los datos crudos en \texttt{k6/}.

\subsection{Instrumentos y herramientas}

\begin{itemize}[nosep]
    \item \textbf{Carga:} k6 v2.2.0 (\texttt{k6/load-test.js})
    \item \textbf{Calidad web:} Lighthouse (vía \texttt{npx lighthouse}), contra build de producción real
    \item \textbf{Seguridad estática:} SpotBugs 4.8.6.4 + find-sec-bugs 1.13.0
    \item \textbf{Seguridad dinámica:} OWASP ZAP baseline scan
    \item \textbf{Cobertura:} JaCoCo 0.8.11
    \item \textbf{Análisis estadístico:} Python 3.14, \texttt{scipy.stats} (Mann-Whitney U), \texttt{numpy}
    \item \textbf{Entorno:} versiones exactas documentadas en \texttt{docs/entorno/versions.txt}
\end{itemize}
\newpage
